% experimental_setup.tex -- Section 3: Experimental Setup

\section{Experimental Setup}
\label{sec:experimental_setup}

We design a controlled ablation study to test whether vision token reduction can preserve closed-loop VLA performance under any configuration. The study varies three independent axes (scoring function, retention ratio, pruning mechanism) around a shared baseline, using a lightweight learned scorer whose training does not alter the pretrained VLA weights. This section describes the scorer, the ablation design, and the evaluation protocol.

\subsection{ACIS Importance Scorer}
\label{sec:acis_scorer}

To select which vision tokens to prune, we train an Attention-weighted Cumulative Importance Score (ACIS) module that predicts a scalar importance score for each of the 256 projected vision tokens (Section~\ref{sec:smolvla_arch}). The ACIS scorer is a lightweight MLP with 3.46M parameters, constituting 0.77\% of SmolVLA's 451.5M total. At inference time, the scorer adds 0.24\% FLOPs overhead and 0.63\% wall-clock latency overhead relative to the unmodified forward pass.

Training follows a Stage-1-only protocol: the ACIS scorer parameters are optimized while all VLA expert weights (vision encoder, MLP connector, LLM backbone, and action head) remain frozen. This separation is deliberate. Fine-tuning the full model during scorer training would conflate two effects: the model could learn to compensate for missing tokens rather than learning to score token importance. By freezing the expert, we ensure that the scorer must identify genuinely important tokens under the original model's representation, and that any subsequent failure cannot be attributed to training-induced degradation of action prediction.

The scorer is trained with an importance margin loss that encourages a clear separation between the scores of retained tokens and those of pruned tokens. Given a retention ratio $k$, the loss penalizes configurations where the top-$k$ and bottom-$(1{-}k)$ token scores overlap, producing a well-separated ranking at convergence. At inference, tokens are sorted by ACIS score and the top-$k$ fraction is kept.

\subsection{Ablation Design}
\label{sec:ablation_design}

The ablation study varies three axes, each designed to rule out a specific failure hypothesis. Table~\ref{tab:ablation_matrix} in Section~\ref{sec:results} reports results for all conditions.

\paragraph{Axis 1: Scoring function.}
We compare the learned ACIS scorer against uniform random token selection. If both produce the same outcome, the failure is not caused by the scorer selecting the wrong tokens.

\paragraph{Axis 2: Retention ratio $k$.}
We test $k \in \{0.5, 0.7, 0.9\}$, corresponding to removing 50\%, 30\%, and 10\% of vision tokens respectively. If even $k{=}0.9$ (removing only 26 of 256 tokens) fails, the collapse is not caused by excessive information loss.

\paragraph{Axis 3: Pruning mechanism.}
We implement three distinct mechanisms for handling pruned tokens:
\begin{itemize}[nosep,leftmargin=*]
  \item \textbf{Zero-mask}: pruned token embeddings are replaced with zero vectors. The sequence length is unchanged, but the model receives out-of-distribution (OOD) activation patterns at the masked positions.
  \item \textbf{True removal}: pruned tokens are physically removed from the input sequence. The LLM processes a shorter sequence with no OOD artifacts, but positional encodings shift for all tokens after the removal point.
  \item \textbf{Token merge}: following the bipartite matching strategy of ToMe~\cite{bolya2023tome}, each pruned token is merged with its nearest retained neighbor by averaging their embeddings. Sequence length and total information content are both preserved.
\end{itemize}
If all three mechanisms fail identically, the collapse is not an artifact of any particular masking strategy.

The three axes are partially crossed. We run ACIS with zero-masking at all three retention ratios ($k{=}0.5, 0.7, 0.9$), random scoring with zero-masking at $k{=}0.5$, and ACIS at $k{=}0.5$ with true removal and merge. This yields six pruning conditions plus the unpruned baseline, covering the critical contrasts needed to isolate each failure hypothesis.

\subsection{Evaluation Protocol}
\label{sec:eval_protocol}

All conditions are evaluated on the three LIBERO suites described in Section~\ref{sec:libero_bench}: Object, Spatial, and Goal. Each suite contains 10 tasks, and we run 50 closed-loop episodes per task, producing 500 episodes per suite per condition.

The baseline condition uses \emph{bypass mode}: the same trained model checkpoint with the ACIS scorer present but pruning disabled, so all 256 vision tokens pass through to the LLM. This ensures that any performance difference between the baseline and pruning conditions is caused by the token reduction itself, not by differences in model weights or the presence of the scorer module.

We report Success Rate (SR) as the primary metric, defined as the fraction of episodes in which the robot completes the full task within the time limit. Because SR is a binary outcome measured over independent episodes, we compute Wilson score 95\% confidence intervals for all reported values. For reference, Wilson 95\% CI for 0 successes out of 500 episodes is $[0.0\%, 0.76\%]$, and for 224 out of 500 is $[40.5\%, 49.2\%]$. These intervals account for the discrete, bounded nature of binary proportions and provide tighter coverage than normal approximation intervals at extreme success rates.
